\section{Conclusion and Future work}
In this work, we introduce and release MiniMax-M1, the world's first open-weight, large-scale reasoning model featuring a lightning attention mechanism. This efficient attention design enables MiniMax-M1 to natively support inputs of up to 1M tokens and generation lengths of 80K tokens---both significantly exceeding capabilities of other open-weight models. These capabilities render MiniMax-M1 uniquely suited for complex, realistic scenarios requiring long context and extended reasoning, properties empirically validated by its strong performance on software engineering, agentic tool use, and long-context understanding benchmarks.
Beyond the inherent efficiency advantages of lightning attention for RL training, this work contributes a novel RL algorithm, \method{}, to accelerate training. Combining architectural advantages with \method{}, we efficiently trained MiniMax-M1, with complete RL training completed in three weeks using 512 H800 GPUs. Across comprehensive evaluations, MiniMax-M1 ranks among the world's best open-weight models alongside DeepSeek-R1 and Qwen3-235B.

Looking forward, as test-time compute continuously scales to power increasingly complex scenarios, we foresee significant potential for such efficient architectures in addressing real-world challenges. These include automating company workflows~\citep{xu2025theagentcompanybenchmarkingllmagents} and conducting scientific research~\citep{si2024llmsgeneratenovelresearch,openai2025research}. Real-world applications particularly demand LRMs that function as agents interacting with environments, tools, computers, or other agents---requiring reasoning across dozens to hundreds of turns while integrating long-context information from diverse sources. We envision MiniMax-M1 serving as a strong foundation for such applications with unique advantages, and we are fully dedicated to further evolving MiniMax-M1 toward this goal.